% SPDX-License-Identifier: GPL-3.0-or-later
% Copyright (C) 2026 Alessandro Greco
\chapter{Корреляция трафика, распознавание сайтов и атаки Sybil}\label{cap:26-attacchi}
\aggiornatoal{23 сентября 2026 г.}{работы, цитируемые в этой главе; tor 0.4.9.13; лабораторная сеть}

\begin{obiettivi}
К концу этой главы вы сможете:
\begin{itemize}[nosep]
  \item классифицировать атаки против анонимности Tor по позиции наблюдения атакующего;
  \item объяснять сквозную корреляцию трафика и измерять её на собственном трафике в лабораторной сети (Лабораторная работа~16);
  \item описывать распознавание сайтов по характеру трафика и ограничения его экспериментальных оценок;
  \item распознавать Sybil-атаки и защитные механизмы сети — от сторожевых узлов до измерения пропускной способности.
\end{itemize}
\end{obiettivi}

Модель угроз из \cref{cap:04-modello-minaccia} ясно показывала, что Tor не рассчитан на противника, наблюдающего оба конца коммуникации. Двадцать лет исследований количественно подтвердили это предупреждение и показали, что точки наблюдения многочисленнее и доступнее, чем может казаться сначала. В этой главе мы рассматриваем атаки с позиции защитника: что требуется атакующему, чего он добивается и какие уже изученные механизмы протокола усложняют задачу.

% ---------------------------------------------------------------------------
\section{Позиция атакующего}\label{sec:att-posizioni}

Полезно классифицировать атаки по тому, что может наблюдать противник:
\begin{description}
  \item[Один конец] Соединение клиент--сторожевой узел, наблюдаемое провайдером связи, локальной сетью или цензором; либо соединение выходной узел--назначение. На линии клиент--сторожевой узел наблюдаются время, объём и направление зашифрованного трафика; между выходом и назначением шифрование прикладного уровня зависит от протокола. Распознавание сайтов может использовать метаданные на стороне клиента.
  \item[Оба конца] Атакующий наблюдает и вход, и выход, либо управляя ретрансляторами в обеих позициях, либо контролируя сети на путях между ними. Это позволяет выполнять корреляцию.
  \item[Внутри сети] Враждебные ретрансляторы, включая промежуточные узлы, могут измерять, задерживать, маркировать или отбрасывать ячейки.
  \item[Над сетью] Базовая инфраструктура Интернета, включая автономные системы (AS) и точки обмена интернет-трафиком (IXP), через которые проходят соединения ретрансляторов.
\end{description}

% ---------------------------------------------------------------------------
\section{Сквозная корреляция}\label{sec:att-correlazione}

Tor намеренно не задерживает и не перемешивает ячейки так, как это делают высоколатентные mix-сети, чтобы оставаться пригодным для интерактивного использования (\cref{sec:radici-latenza}). Поэтому \emph{форма} трафика — сколько байт проходит в каждом временном интервале — часто остаётся узнаваемой с двух концов, несмотря на задержки, буферизацию, мультиплексирование и посторонний трафик (\cref{fig:correlazione}). Тот, кто наблюдает оба конца, может сравнивать формы и решать, соответствуют ли потоки друг другу, не взламывая шифрование.

\begin{figure}[htbp]
\centering
% SPDX-License-Identifier: GPL-3.0-or-later
% Copyright (C) 2026 Alessandro Greco
%% Русское альтернативное описание EPUB.
\altfig{Синтетическая иллюстрация, а не реальные измерения, корреляции трафика Tor в одном направлении. Верхний график показывает относительный объём зашифрованного трафика между клиентом и guard за интервалы по 250 миллисекунд; нижний — относительный объём между exit и сервером. Два независимо заданных условных ряда содержат похожие, но не идентичные всплески, сдвинутые в примере на один интервал; их высота различается и добавлен шум. Шифрование не обязательно скрывает временные и объёмные метаданные. Буферизация, накладные расходы, мультиплексирование и посторонний трафик могут менять профили. Фиксированная задержка не является свойством Tor, а статистическое сопоставление не всегда верно.}
% Синтетические ряды, не наблюдения
\begin{tikzpicture}[font=\sffamily\scriptsize,x=0.32cm,y=0.09cm]
  \def\serieIngresso{0/6,1/0,2/0,3/9,4/12,5/2,6/0,7/0,8/0,9/5,10/14,11/14,12/3,13/0,14/0,15/7,16/0,17/0,18/11,19/4,20/0,21/0,22/8,23/2}
  \def\serieUscita{0/0,1/5,2/0,3/0,4/8,5/11,6/3,7/0,8/0,9/1,10/4,11/13,12/12,13/4,14/0,15/0,16/6,17/0,18/1,19/10,20/5,21/0,22/1,23/7,24/3}
  \foreach \x/\h in \serieIngresso {
    \fill[torviola!70] (\x,20) rectangle ++(0.8,\h);
  }
  \foreach \x/\h in \serieUscita {
    \fill[torverde!70!black] (\x,0) rectangle ++(0.8,\h);
  }
  \draw[->] (-0.5,20) -- (27,20) node[right] {$t$};
  \draw[->] (-0.5,0) -- (27,0) node[right] {$t$};
  \node[left,align=right] at (-0.6,26) {наблюдатель\\на входе\\(клиент--guard)};
  \node[left,align=right] at (-0.6,6) {наблюдатель\\на выходе\\(exit--сервер)};
  \draw[<->,torgrigio] (10,17) -- node[above,font=\sffamily\tiny] {пример: сдвиг на один интервал} (11,17);
  \node[text=torgrigio] at (13,40) {относительный объём за 250\,мс};
  \node[text=torgrigio,font=\sffamily\tiny] at (13,36.5) {синтетический пример, не измерения};
\end{tikzpicture}

\caption{Два \emph{синтетических} ряда, а не реальные измерения объёма трафика в одном направлении на линиях клиент--guard и exit--сервер. Всплески похожи, но не идентичны; сдвиг на один интервал 250\,мс выбран лишь для иллюстрации. Шифрование скрывает содержимое, но не всегда временные и объёмные признаки. Буферизация, накладные расходы, мультиплексирование и дополнительные потоки изменяют картину; статистическое сходство само по себе не доказывает совпадение.}
\label{fig:correlazione}
\end{figure}

\paragraph{Насколько легко сопоставлять потоки?} В 2018 году DeepCorr обучил нейронную сеть распознавать пары Tor-потоков. При примерно 900 пакетах на поток он достиг 96\,\% точности против 4\,\% у более раннего статистического метода в тех же условиях~\cite{nasr2018deepcorr}. Для атакующего основная трудность обычно заключается не в сравнении потоков, а в занятии правильных точек наблюдения.

\paragraph{Насколько вероятно наблюдать оба конца?} Johnson и коллеги месяцами симулировали типичных пользователей с реальным алгоритмом выбора маршрута Tor~\cite{johnson2013users}. При вредоносных ретрансляторах умеренной пропускной способности 80\,\% пользователей во всех смоделированных типах были деанонимизированы хотя бы один раз в течение шести месяцев. Для одного AS пользователи в некоторых распространённых географических положениях приближались к 100\,\% в течение трёх месяцев. Эти исторические результаты симуляции не означают, что каждое соединение раскрыто; они показывают, что возможности накапливаются со временем. Именно такая логика стала одной из причин перехода к более долгоживущим сторожевым узлам и меньшему числу сторожевых узлов на клиента (\cref{sec:guard}).

\paragraph{Интернет как противник.} RAPTOR описал три способа, которыми AS может увеличить свои возможности~\cite{sun2015raptor}: использовать асимметричные интернет-маршруты, поскольку иногда достаточно наблюдать \emph{одно} направление на каждом конце, включая подтверждения TCP; использовать естественные изменения маршрутизации со временем; а также перехват или перенаправление префиксов посредством BGP. Эти методы не меняют Tor. Защита включает мониторинг маршрутизации и потенциально выбор путей с учётом топологии.

\paragraph{Активные атаки.} Атакующий, контролирующий один конец, может намеренно \emph{впечатывать} отличимый трафиковый рисунок, задерживая или группируя пакеты, а затем искать его на другом конце. Атака маркировки, обсуждавшаяся на уровне шифрования в \cref{sec:tor1-debolezze}, — крайняя форма этой идеи. Уже в 2005 году Murdoch и Danezis показали, что наблюдать реальные каналы не всегда обязательно. Вредоносный сервер навязывает ритм трафику к жертве; атакующий измеряет задержку на собственных цепях через ретрансляторы-кандидаты, и те, чья задержка меняется в том же ритме, могут находиться в цепи жертвы~\cite{murdoch2005lowcost}.

% ---------------------------------------------------------------------------
\section{Лабораторная работа: корреляция собственного трафика}\label{sec:lab-attacchi}

\begin{laboratorio}[16]{Сквозная корреляция}
\textbf{Цель.} Показать в лабораторной сети, что рисунок трафика сохраняется через цепь и позволяет определить, какой поток на выходной стороне соответствует какому клиенту.\\
\textbf{Предварительные условия.} Лабораторная работа~11; лабораторная сеть должна быть запущена. Все потоки создаются самим экспериментальным сценарием.\\
\textbf{Оценочное время.} 1 час.
\end{laboratorio}

Скрипт \file{labs/ru/лаборатория/correlation.py} запускает сервер на порту контейнера 8095 и шесть клиентов. Каждый клиент использует разные SOCKS-учётные данные и при включённой опции \opt{IsolateSOCKSAuth} их потоки не разделяют одну цепь (\cref{sec:lab-stem}). В течение 20 секунд клиенты случайным образом чередуют серии блоков по 2\,KB с паузами. Два \enquote{наблюдателя} записывают данные:
\begin{itemize}
  \item Входной наблюдатель записывает моменты, когда каждый клиент отправляет данные. Это идеализирует картину наблюдателя на линии к сторожевому узлу. Там данные были бы упакованы в ячейки и TLS, но нужный для корреляции временной рисунок переживает оба слоя.
  \item Наблюдатель на выходной стороне, наш сервер, записывает моменты поступления байтов по каждому соединению, и получает четырёхбайтовый идентификатор для проверки результатов, который не участвует в выборе соответствия.
\end{itemize}
Оба агрегируют байты в интервалы по 250\,мс. Для каждого клиента вычисляется коэффициент корреляции Пирсона со всеми шестью потоками на стороне exit. Наиболее коррелированный поток считается предполагаемым соответствием. Скрипт знает правильный ответ, передаваемый клиентом в первых четырёх байтах, но использует его только для оценки, а не для выбора соответствия.

\begin{lstlisting}[language=bash,numbers=none]
$ docker exec -w /tmp torbook-network python3 /lab/correlation.py 1
клиент 0: +0.79 -0.11 +0.12 +0.17 +0.11 -0.01  -> поток exit 0 ВЕРНО
клиент 1: -0.11 +0.80 -0.02 +0.07 +0.18 +0.08  -> поток exit 1 ВЕРНО
клиент 4: +0.13 +0.02 +1.00 +0.21 +0.10 +0.21  -> поток exit 2 ВЕРНО
клиент 3: +0.21 +0.09 +0.33 +0.81 -0.07 +0.03  -> поток exit 3 ВЕРНО
клиент 2: +0.17 +0.05 +0.10 -0.09 +0.85 -0.03  -> поток exit 4 ВЕРНО
клиент 5: -0.00 +0.01 +0.23 +0.06 -0.13 +0.92  -> поток exit 5 ВЕРНО
правильных совпадений: 6/6
\end{lstlisting}
Столбцы обозначают выходные потоки в порядке установления соединений, который отличается от порядка клиентов. В этом иллюстративном выводе правильные коэффициенты лежат между 0,79 и 1,00; порядок потоков и значения могут меняться от запуска к запуску. Не нужны ни ключи расшифрования, ни расшифрованные ячейки — только объём данных и время. Лабораторная сеть мала и почти не нагружена, поэтому задержки минимальны. Исторические результаты DeepCorr показывают, что значительно более сложные классификаторы могут работать с несколькими сотнями пакетов в реальной сети при исследованных условиях. Это фрагмент эталонного запуска, а не точный прогноз: возможны и 5/6 совпадений. Лабораторная работа входит в \texttt{make verifica-lab}.

% ---------------------------------------------------------------------------
\section{Распознавание сайтов по трафику}\label{sec:att-wf}

Атакующий, наблюдающий только линию клиент--сторожевой узел, не может выполнить сквозную корреляцию, но может попытаться \emph{распознать} посещаемый сайт. Каждая веб-страница порождает характерную последовательность ячеек в обоих направлениях: количество, порядок и серии передач. Атакующий посещает целевые сайты через Tor, записывает их последовательности и обучает классификатор. Затем он сравнивает трафик жертвы с изученными паттернами.

Исследования атак много лет представляют собой соревнование между классификаторами и защитой:
\begin{itemize}
  \item В 2011 году Panchenko и коллеги использовали объём, время и направление трафика, повысив распознавание в закрытом наборе из 775 сайтов с 3\,\% до 55\,\%. Они также продемонстрировали эффективную атаку в \emph{открытом мире}, где жертва может посещать сайты вне обучающего набора атакующего~\cite{panchenko2011website}.
  \item В 2018 году Deep Fingerprinting использовал свёрточную нейронную сеть и превысил 98\,\% точности на 95 незащищённых сайтах и 90\,\% против WTF-PAD — защиты на основе адаптивного заполнения, рассматривавшейся для Tor~\cite{sirinam2018deep}.
\end{itemize}

\begin{approfondimento}[Ошибка базовой частоты]
Очень высокие значения точности требуют осторожной интерпретации. Juarez и коллеги показали, что многие исследования предполагали выгодные условия: посещается одна страница за раз, у жертвы и атакующего одна версия Tor Browser, а обучающие данные собраны незадолго до атаки~\cite{juarez2014critical}. Ещё важнее то, что в открытом мире отслеживаемые сайты составляют крошечную долю всех посещений. Даже классификатор с низкой частотой ложных срабатываний может выдавать больше ложных тревог, чем правильных распознаваний, среди миллионов нерелевантных посещений. Это та же статистическая ловушка, что и при скрининговых тестах редких заболеваний; реальному атакующему приходится учитывать её.
\end{approfondimento}

\paragraph{Защита.} Заполнение канала (\cref{cap:07-canali}) затрудняет определение периодов бездействия, но не предназначен специально для защиты от распознавания страниц по трафику. tor также реализует более общий механизм \emph{заполнения цепей}: конечные автоматы могут планировать фиктивные ячейки. На эталонную дату исходного текста развёрнутые машины помогают цепям представления и встречи onion-сервисов выглядеть как обычные цепи (\cref{cap:14-onion-services}, \proposta{302}). Эффективная широкая защита от распознавания сайтов по трафику с приемлемыми накладными расходами по пропускной способности остаётся открытой исследовательской задачей.

% ---------------------------------------------------------------------------
\section{Sybil: размножение идентичностей ретрансляторов}\label{sec:att-sybil}

\textbf{Sybil-атака} вводит в сеть множество идентичностей. В Tor атакующий добавляет многочисленные ретрансляторы в надежде чаще занимать позиции сторожевых и выходных узлов. Три защитных механизма повышают стоимость:
\begin{description}
  \item[Измеренная пропускная способность] Ретранслятор, завышающий заявленную пропускную способность, мог бы привлечь слишком много клиентов. Bauer и коллеги продемонстрировали такую атаку через дескрипторы в 2007 году~\cite{bauer2007lowresource}. Сегодня веса консенсуса основаны на измерениях, учитываемых центрами каталогов, а не на одних заявлениях (\cref{sec:dir-voto}); чтобы привлечь значительный трафик, нужно предоставить реальную пропускную способность.
  \item[Ограничения маршрута] Два ретранслятора из одного семейства или одной IPv4-сети /16 не могут находиться в одной цепи (\cref{sec:percorso-vincoli}). Честные операторы объявляют семейства; нечестных хотя бы ограничивают адреса.
  \item[Мониторинг сети] Центры каталогов могут исключать ретрансляторы или присваивать \texttt{BadExit}. Добровольцы и исследователи изучают группы с общими характеристиками и наблюдают поведение выходных узлов. Например, в 2014 году сканер выходных узлов обнаружил 25 ретрансляторов, изменявших трафик или имевших опасную конфигурацию, включая перехват HTTPS, изменение DNS и \emph{sslstrip}~\cite{winter2014spoiled}.
\end{description}
В \cref{cap:27-casi-rete} рассматриваются вредоносные группы ретрансляторов, найденные в последующие годы, и показано, почему мониторинг остаётся постоянной задачей.

\paragraph{Sybil-атаки против onion-сервисов.} В v2 позиции HSDir, ответственных за сервис, были предсказуемы. Атакующий мог генерировать идентичности ретрансляторов так, чтобы они занимали нужные позиции. Biryukov и коллеги в 2013 году использовали это для измерения популярности сервисов, создания недоступности и облегчения атак на местоположение~\cite{biryukov2013trawling}. В v3 позиции зависят от ослеплённых ключей и общего случайного значения, меняющегося ежедневно (\cref{cap:14-onion-services}), что сильно усложняет заранее рассчитанное размещение; после публикации случайного значения текущего периода ответственные HSDir всё ещё можно вычислить.

% ---------------------------------------------------------------------------
\section{Обнаружение сторожевого узла}\label{sec:att-guard}

Сторожевой узел защищает клиента, но одновременно является привлекательной целью: его идентификация показывает атакующему, где наблюдать. Против onion-сервиса атакующий может провоцировать множество построений цепей, например создавая много соединений, пока вредоносный ретранслятор не окажется промежуточным узлом непосредственно рядом со сторожевым. Первая атака такого типа, опубликованная в 2006 году, сама стала одним из стимулов к использованию сторожевых узлов~\cite{overlier2006locating}. Vanguards (\cref{sec:vanguards}) добавляют слой промежуточных ретрансляторов, фиксируемых на некоторое время. Это затрудняет продвижение атакующего к настоящему сторожевому узлу, но не гарантирует невозможность его обнаружения.

% ---------------------------------------------------------------------------
\section*{Итоги}
\begin{itemize}
  \item Атаки различаются по точке наблюдения: один конец, оба конца, внутри сети или базовые интернет-маршруты.
  \item Корреляция сравнивает форму трафика на обоих концах без взлома шифрования. Главная задача атакующего — занять нужные позиции; возможности накапливаются со временем.
  \item Распознавание сайтов по трафику определяет страницы по зашифрованному трафику на одном конце. Оценки должны учитывать базовые частоты и реалистичные условия.
  \item Защита от Sybil включает измеренные веса, семейства, ограничения /16 и мониторинг выходных узлов и групп ретрансляторов. В v3 размещением HSDir сложнее манипулировать заранее, хотя узлы, ответственные за текущий период, можно вычислить.
  \item Сторожевые узлы и vanguards делают обнаружение первого перехода более медленным и дорогостоящим.
\end{itemize}

% ---------------------------------------------------------------------------
\section*{Упражнения}

\begin{esercizio}[1]\label{es:26-posizione}
Для каждой атаки в этой главе укажите минимальную позицию наблюдения атакующего и защиту Tor, которая усложняет атаку.
\end{esercizio}

\begin{esercizio}[1]\label{es:26-base}
Классификатор распознаёт отслеживаемый сайт с вероятностью 0,95 и ошибочно помечает нерелевантное посещение с вероятностью 0,001. Если только одно посещение из 100\,000 действительно относится к отслеживаемому сайту, какая доля срабатываний будет правильной?
\end{esercizio}

\begin{esercizio}[2]\label{es:26-bin}
Повторите Лабораторную работу~16 с двухсекундными интервалами, затем с длительностью всего пять секунд и, наконец, заставив всех клиентов использовать одинаковый ритм трафика. Какое изменение сильнее всего ухудшает сопоставление и почему?
\end{esercizio}

\begin{esercizio}[2]\label{es:26-difesa}
Измените клиентов Лабораторной работы~16 так, чтобы они отправляли блоки с постоянной скоростью, отправляя фиктивные блоки при отсутствии данных. Каковы дополнительные затраты пропускной способности и что происходит с корреляцией?
\end{esercizio}

\begin{esercizio}[3]\label{es:26-probabilita}
Используя \file{path\_selection.py} и реальный консенсус из Лабораторной работы~15, вычислите вероятность того, что атакующий, контролирующий 5\,\% пропускной способности сторожевых узлов и 5\,\% пропускной способности выходных узлов, займёт оба конца хотя бы одной клиентской цепи в течение шести месяцев. Сравните фиксированный сторожевой узел с выбором нового для каждой цепи. Задайте условное число цепей за шесть месяцев; лишь для учебной модели считайте выбор независимым и взвешенным по пропускной способности, пояснив пределы такого расчёта.
\end{esercizio}

% ---------------------------------------------------------------------------
\section*{Дополнительная литература}
\begin{itemize}
  \item A.~Johnson et al., \titolo{Users Get Routed}~\cite{johnson2013users}.
  \item Y.~Sun et al., \titolo{RAPTOR: Routing Attacks on Privacy in Tor}~\cite{sun2015raptor}.
  \item M.~Juarez et al., \titolo{A Critical Evaluation of Website Fingerprinting Attacks}~\cite{juarez2014critical}.
  \item \emph{Anonymity Bibliography}, \url{https://www.freehaven.net/anonbib/}, для более поздних исследований.
\end{itemize}
